\textbf{模意义组合数选法}

\textbf{输入与边界}
求 $\binom nk\bmod M$，其中 $n,k\ge0$、$M\ge1$。
$k>n$ 时结果为 $0$；$M=1$ 时结果为 $0$。
以下默认 $0\le k\le n$，可先令 $k\leftarrow\min(k,n-k)$。

\textbf{模素数 $p$}
\begin{itemize}
    \item 若 $n<p$ 且要回答多次查询：预处理 $0,\ldots,n$ 的阶乘和逆阶乘，
    单次 $\binom nk=n!/(k!(n-k)!)\pmod p$ 可在 $O(1)$ 求得。
    \item 若 $n\ge p$：使用 Lucas 定理。写
    $n=\sum_i n_ip^i$、$k=\sum_i k_ip^i$，则
    $\binom nk\equiv\prod_i\binom{n_i}{k_i}\pmod p$；
    某位 $k_i>n_i$ 即为 $0$。
    \item 若 $k<p$ 且 $k$ 很小，也可直接计算
    $\prod_{j=0}^{k-1}(n-j)/(k-j)\pmod p$；
    分母均与 $p$ 互质。一般 $k\ge p$ 时不能这样逐项取逆。
\end{itemize}
Lucas 的位数为 $O(\log_p n)$；每位组合数可用长度 $p$ 的阶乘表，
或按较小下指标逐项相乘。若 $p$ 很大，预处理整张表未必可行。

\textbf{模合数 $M$}
分解 $M=\prod_i p_i^{e_i}$，先分别计算
$\binom nk\bmod p_i^{e_i}$，再用中国剩余定理合并。
计算模 $p^e$ 时先求
\[
v_p\!\left(\binom nk\right)
=\sum_{j\ge1}\left(
\left\lfloor\frac n{p^j}\right\rfloor
-\left\lfloor\frac k{p^j}\right\rfloor
-\left\lfloor\frac{n-k}{p^j}\right\rfloor\right).
\]
若此值 $\ge e$，结果为 $0$；否则剥离阶乘中的所有 $p$ 因子，
只对剩余的可逆部分求逆，最后乘回 $p^{v_p}$。
可参见 \texttt{数学/数论/ExLucas·扩展Lucas.cpp}；
该实现按各质数幂大小预处理，质数幂很大时不适合直接套用。

\textbf{常见误用}
模合数时，即使 $n<M$，$k!$ 也可能没有模逆元；
模素数但 $n\ge p$ 时，直接套阶乘逆元公式同样会遇到 $n!\equiv0\pmod p$。
先检查分母是否可逆，再决定是否使用阶乘公式。
